\begin{table}[H]
\caption{Carteira Selecionada contra Referências — CS:GO}\label{tab:portfolio_benchmarks} 

\fontsize{12.0pt}{14.4pt}\selectfont

\begin{tabular*}{\linewidth}{@{\extracolsep{\fill}}lrrrr}
\toprule
\textbf{Carteira} & \textbf{ROI líquido} & \textbf{P5} & \textbf{P95} & \textbf{Percentil da selecionada} \\
\midrule\addlinespace[2.5pt]
Carteira selecionada (k = 10) & 36,42\% & n/d & n/d & n/d \\
Todos os candidatos & 46,05\% & n/d & n/d & n/d \\
Todos os elegíveis & 53,86\% & n/d & n/d & n/d \\
Item mediano entre os candidatos & 35,23\% & n/d & n/d & n/d \\
Aleatória entre elegíveis (média) & 53,66\% & 25,19\% & 98,68\% & 22\% \\
Aleatória entre candidatos (média) & 46,30\% & 21,63\% & 82,26\% & 32\% \\
\bottomrule
\end{tabular*}
\par\vspace{\belowcaptionskip}\footnotesize Fonte: Elaboração própria.\par
\begin{minipage}{\linewidth}
Compra e retenção com pesos iguais e taxa de venda de 15\%. As carteiras aleatórias têm 10 itens e vêm de 1.000 sorteios. Percentil: parcela dos sorteios com ROI igual ou menor ao da selecionada.\\
\end{minipage}
\end{table}
